\documentclass{article}
\usepackage[margin=2cm]{geometry}
\usepackage[ngerman]{babel}
\usepackage{float}
\usepackage{graphicx}
\usepackage{comment}
\usepackage[version=4]{mhchem}

\title{OC3 Skript}
\author{}
\date{\today}

\begin{document}

\maketitle
\tableofcontents
\newpage

\section{OC3}
\subsection{Aromatizität}
Die\textbf{ moderne Definition} der Aromatizität besagt: Eine Verbinfung ist armoatisch wenn alle Ringatome in ein einziges konjungiertes System 
einbezogen sind und diese einen dynamischen Ringstrom besitzen.\\

\noindent Zudem gilt die altbekannte \textbf{ Hückel-Regel}:
\begin{enumerate}
	\item monocyclische, vollständig konjungierte, planare Syseme
	\item $4n+2=\pi Elektronen$
\end{enumerate}
Abseits hiervon gelten folgende Kriterien: 
\begin{enumerate}
	\item Es herrscht Bindungslängen-Ausgleich (Alternanz)
\end{enumerate}
\textbf{Aromatische Systeme }sind stabiler als ihre acyclischen Analoga.\\
\textbf{anti-aromatische Systeme} sind weniger stabil als ihre cyclischen Analoga. \\
\textbf{nicht aromatische Systeme} besitzen ähnliche Delokalisierungsenthalpien wie ihre acyclischen Analoga.

\subsection{Annulene}
\subsection{Polycyclische Aromaten}
Die Verbindungen der polycyclischen aromatischen Kohlenwasserstoffe können in zwei Kaegorien unterteilt werden. Einmal den nicht kodensierten und den kondensierten Aromaten. Ein kondensierter Aromat ist hierbei eine Verbindung bei der zwei oder mehr Ringe über mehr als eine Bindung miteinander verbunden sind. Des weiteren lassen sich die kondensierten 
Aromaten in zwei weitere Kategorien unterteilen. Es existieren benzoide und nicht bezoide Verbindungen. Eine benzoide Verbindung ist eine Verbindung, 
die aus mehr mehrern kondensierten Benzolringen besteht. Benzoide Verbindungen bestehen aus Ringen einer anderen Ringgröße. Sowohl benzoide als auch nict benzoide Verbindungen sind aromatisch. \\

Für benzoide Aromaten sind folgende Regeln definiert: 
\begin{enumerate}
	\item Die Resonazstabilisierung nimmt mit der Anzahl an Ringen ab.
	\item Je größer die Anzahl an Sextetten in einem Aromaten desto stabiler ist diese.
\end{enumerate}
\subsubsection{linear anellierte polycyclische Verbindungen}
Hiebei handelt es sich um eine gerade lineare Reihe von kodensierten Benzolringen.
\subsubsection{nicht-linear kondensierte benzoide Systeme}
\subsubsection{nicht benzoide polycyclische Systeme}
\subsubsection{Kreuzkomjugierte Systeme}

\subsection{Heterocyclen}
\subsubsection{Hantzsch-Widman-Nomenklatur}
\subsubsection{Pyrrol} 
\subsubsection{Porphyrine}
\subsubsection{Imidazol}
\subsubsection{Indol}
\subsubsection{Pyridin}
\subsubsection{Furan}
\subsubsection{Benzofuran}
\subsubsection{Isobenzofuran}

\subsection{Pericyclische Reaktionen}
\subsubsection{Cycloadditionen}
\subsubsection*{1,3-dipolare Cycloaddition}
\subsubsection{Sigmatrope Umlagerung}
\subsubsection{Elektrocyclische Reaktionen}
\subsubsection{En-Reaktion}
\subsubsection{Chelotrope Reaktionen}

\end{document}